% =============================================================================
% CAPÍTULO 5 — RESULTADOS Y EVALUACIÓN
% =============================================================================
\chapter{Resultados y Evaluación}
\label{ch:resultados}

El presente capítulo reporta los resultados obtenidos con \textsc{macau} en el
contexto de la carrera \programaTesis\ de la Universidad Católica de Salta.
Se organiza en dos grandes bloques: primero se presenta el prototipo funcional
mediante capturas del sistema en operación, para dar cuenta de la solución
construida antes de entrar en la evaluación cuantitativa; luego se reportan
las métricas experimentales de los módulos de extracción automática, validación
normativa y asistencia IA.

%%────────────────────────────────────────────────────────────────────────────
\section{Presentación del prototipo funcional}
\label{sec:prototipo}
%%────────────────────────────────────────────────────────────────────────────

Esta sección presenta el sistema \textsc{macau} en funcionamiento a través de
capturas organizadas por flujo de usuario. Se documentan con detalle los
recorridos de autenticación, recuperación de contraseña, y los paneles
propios de los roles Directivo y Docente, que concentran las funcionalidades
centrales del sistema. Las vistas de Secretario y Miembro de la Comisión
Curricular se documentan en el Anexo~\ref{anexo:capturas}.

\subsection{Autenticación}
\label{subsec:cap_autenticacion}

La pantalla de inicio de sesión es el punto de entrada único para todos los
roles del sistema (Figura~\ref{fig:cap_login}). El formulario solicita correo
electrónico y contraseña; ante credenciales inválidas se muestra un mensaje
de error genérico sin revelar si el correo existe en el sistema, lo que
mitiga la enumeración de usuarios (Figura~\ref{fig:cap_login_error}).
Un inicio exitoso redirige al dashboard del rol asignado o, en el caso de
cuentas multi-rol, al selector de contexto descrito en la sección siguiente.

\begin{figure}[H]
  \centering
  \fbox{\begin{minipage}[c][5.5cm][c]{0.85\linewidth}
    \centering{\color{gray}\textit{[Captura pendiente: pantalla de inicio de sesión — formulario vacío]}}
  \end{minipage}}
  \caption[Pantalla de inicio de sesión de \textsc{macau}.]{Pantalla de inicio
    de sesión. Se solicitan correo institucional y contraseña; no existe
    registro público de cuentas de usuario.}
  \label{fig:cap_login}
\end{figure}

\begin{figure}[H]
  \centering
  \fbox{\begin{minipage}[c][5.5cm][c]{0.85\linewidth}
    \centering{\color{gray}\textit{[Captura pendiente: mensaje de error tras ingresar credenciales incorrectas]}}
  \end{minipage}}
  \caption[Mensaje de error por credenciales incorrectas.]{El sistema informa
    el fallo de autenticación sin especificar cuál de los dos campos es
    incorrecto, reduciendo la superficie de ataque por enumeración.}
  \label{fig:cap_login_error}
\end{figure}

\subsection{Selección de rol en cuentas multi-rol}
\label{subsec:cap_multirrol}

Cuando una cuenta tiene asignados dos o más roles —por ejemplo, un docente
que también integra la Comisión Curricular— el sistema presenta un selector
de rol inmediatamente después de la autenticación exitosa
(Figura~\ref{fig:cap_multirrol}). La sesión activa queda asociada al rol
seleccionado; para cambiar de contexto el usuario debe cerrar sesión y
elegir un rol diferente en el siguiente inicio.

\begin{figure}[H]
  \centering
  \fbox{\begin{minipage}[c][5.5cm][c]{0.85\linewidth}
    \centering{\color{gray}\textit{[Captura pendiente: modal de selección de rol para cuentas con múltiples roles asignados]}}
  \end{minipage}}
  \caption[Selector de rol para usuarios con múltiples roles.]{Pantalla de
    selección de contexto: el usuario elige bajo qué rol desea operar. Cada
    rol habilita un conjunto de vistas y acciones diferente según la tabla
    de permisos RBAC (Tabla~\ref{tab:rbac}).}
  \label{fig:cap_multirrol}
\end{figure}

\subsection{Recuperación de contraseña}
\label{subsec:cap_reset}

El sistema contempla tres escenarios de cambio de credenciales. El primero
corresponde al primer inicio de sesión de una cuenta creada por el Directivo,
donde las credenciales provisorias son correo~/ correo y el sistema obliga
el cambio antes de dar acceso al dashboard
(Figura~\ref{fig:cap_reset_primer_login}). El segundo es la recuperación
iniiciada por el usuario desde la pantalla de login
(Figura~\ref{fig:cap_reset_solicitud}): se ingresa el correo y se recibe
un enlace de un solo uso con vigencia de 24 horas. El tercero es el
establecimiento de la nueva contraseña desde ese enlace
(Figura~\ref{fig:cap_reset_nueva}).

\begin{figure}[H]
  \centering
  \fbox{\begin{minipage}[c][5.5cm][c]{0.85\linewidth}
    \centering{\color{gray}\textit{[Captura pendiente: cambio de contraseña obligatorio en el primer inicio de sesión con credenciales provisorias]}}
  \end{minipage}}
  \caption[Cambio de contraseña obligatorio en el primer inicio de sesión.]{Al
    acceder por primera vez con las credenciales provisorias asignadas por
    el Directivo, el sistema bloquea el acceso al dashboard y redirige
    al formulario de cambio de contraseña.}
  \label{fig:cap_reset_primer_login}
\end{figure}

\begin{figure}[H]
  \centering
  \fbox{\begin{minipage}[c][5.5cm][c]{0.85\linewidth}
    \centering{\color{gray}\textit{[Captura pendiente: formulario de solicitud de recuperación de contraseña desde la pantalla de login]}}
  \end{minipage}}
  \caption[Solicitud de recuperación de contraseña.]{El usuario ingresa su
    correo institucional; el sistema envía un enlace de reseteo de un solo
    uso si la cuenta existe, sin confirmar ni denegar su existencia.}
  \label{fig:cap_reset_solicitud}
\end{figure}

\begin{figure}[H]
  \centering
  \fbox{\begin{minipage}[c][5.5cm][c]{0.85\linewidth}
    \centering{\color{gray}\textit{[Captura pendiente: formulario de nueva contraseña accedido desde el enlace de reseteo recibido por correo]}}
  \end{minipage}}
  \caption[Establecimiento de nueva contraseña.]{Formulario accesible
    únicamente desde el enlace de reseteo; el token queda invalidado al
    ser consumido o al expirar el plazo de 24 horas.}
  \label{fig:cap_reset_nueva}
\end{figure}

\subsection{Vista del Directivo}
\label{subsec:cap_directivo}

El Directivo accede al dashboard central del sistema, desde el cual gestiona
propuestas de todas las carreras habilitadas. La Figura~\ref{fig:cap_dir_dashboard}
muestra el listado de propuestas con indicadores de estado por color,
filtros por carrera, docente y año lectivo, y acceso rápido a las acciones
de cada propuesta. Desde este mismo panel puede vincular carpetas de
Google Drive a cada carrera y acceder al plan de trabajo de acreditación.

\begin{figure}[H]
  \centering
  \fbox{\begin{minipage}[c][7cm][c]{0.88\linewidth}
    \centering{\color{gray}\textit{[Captura pendiente: dashboard del Directivo — listado de propuestas con filtros, estados y acciones rápidas]}}
  \end{minipage}}
  \caption[Dashboard del rol Directivo.]{Panel principal del Directivo:
    tabla de propuestas con carrera, docente, estado semáforo y fecha;
    filtros laterales; menú de acciones sobre cada propuesta.}
  \label{fig:cap_dir_dashboard}
\end{figure}

\begin{figure}[H]
  \centering
  \fbox{\begin{minipage}[c][7cm][c]{0.88\linewidth}
    \centering{\color{gray}\textit{[Captura pendiente: vista de propuesta con reporte de validación normativa (conformidades y no-conformidades con referencia a la norma)]}}
  \end{minipage}}
  \caption[Vista de propuesta con resultado de validación normativa.]{Vista
    de detalle: campos extraídos en el área principal y, en el panel
    lateral, el reporte del agente de validación con conformidades,
    no-conformidades y referencia a la normativa aplicable (Ord.\ 75/23).}
  \label{fig:cap_dir_propuesta}
\end{figure}

\subsection{Vista del Docente}
\label{subsec:cap_docente}

El Docente opera sobre sus propias propuestas. El flujo habitual comienza
con la importación del documento (Figura~\ref{fig:cap_doc_importar}),
que puede realizarse desde un archivo DOCX mediante drag-and-drop o desde
una URL de Google Docs. El proceso dispara la extracción automática de
campos; la Figura~\ref{fig:cap_doc_ficha} muestra la ficha resultante con
todos los campos estructurados listos para revisión y edición inline.
Una vez que el Directivo emite notificaciones sobre la propuesta, estas
aparecen en el panel de notificaciones del docente.

\begin{figure}[H]
  \centering
  \fbox{\begin{minipage}[c][5.5cm][c]{0.88\linewidth}
    \centering{\color{gray}\textit{[Captura pendiente: pantalla de importación — tabs DOCX / Google Docs con zona de carga y campo de URL]}}
  \end{minipage}}
  \caption[Pantalla de importación de propuesta por el Docente.]{Interfaz
    de importación con dos modos: archivo DOCX por arrastrar-soltar,
    o URL de Google Docs para importación directa.}
  \label{fig:cap_doc_importar}
\end{figure}

\begin{figure}[H]
  \centering
  \fbox{\begin{minipage}[c][7cm][c]{0.88\linewidth}
    \centering{\color{gray}\textit{[Captura pendiente: ficha de propuesta con campos extraídos automáticamente: título, fundamentación, RA, contenidos mínimos, bibliografía]}}
  \end{minipage}}
  \caption[Ficha de propuesta con campos extraídos.]{Ficha generada luego
    de la importación: todos los campos estructurados con el contenido
    extraído automáticamente por los agentes, editables inline por el
    docente antes de enviar a revisión.}
  \label{fig:cap_doc_ficha}
\end{figure}

%%────────────────────────────────────────────────────────────────────────────
\section{Diseño del experimento de evaluación}
\label{sec:disenio_experimento}
%% ─────────────────────────────────────────────────────────────

% TODO: Describir el setup experimental:
%   - Corpus usado: N propuestas de la carrera \programaTesis
%   - Cómo se construyó el ground truth (anotación manual, quién anotó,
%     criterios de anotación, acuerdo entre anotadores si aplica)
%   - División train/test si corresponde (o si es evaluación sobre
%     todo el corpus)
%   - Entorno de ejecución: hardware, versiones de software
%   - Reproducibilidad: semillas, modelos fijos
\ldots

%% ─────────────────────────────────────────────────────────────
\section{Resultados de extracción automática de campos}
\label{sec:resultados_extraccion}
%% ─────────────────────────────────────────────────────────────

% TODO: Reportar métricas de extracción por campo:
%   - Tabla: Campo | Precisión | Recall | F1 | N muestras
%     (todos los campos: título, fundamentación, objetivos,
%     contenidos mínimos, RA conceptuales/procedimentales/actitudinales,
%     metodología, evaluación, bibliografía)
%   - Análisis de los campos con mejor/peor performance
%   - Ejemplos concretos de extracción correcta e incorrecta
%   - Casos de falla: qué variantes de formato rompen el extractor
\ldots

%% ─────────────────────────────────────────────────────────────
\section{Resultados de validación normativa automática}
\label{sec:resultados_validacion}
%% ─────────────────────────────────────────────────────────────

% TODO: Reportar el desempeño del agente de validación normativa:
%   - ¿Cuántas no-conformidades reales detectó correctamente?
%   - ¿Cuántas no-conformidades inventó (falsos positivos)?
%   - Tipos de no-conformidad más frecuentes en el corpus
%   - Ejemplo de reporte de no-conformidad generado por el sistema
%   - Comparación con revisión manual (si se realizó)
\ldots

%% ─────────────────────────────────────────────────────────────
\section{Resultados de asistencia LLM}
\label{sec:resultados_llm}
%% ─────────────────────────────────────────────────────────────

% TODO: Evaluar la calidad de las sugerencias generadas por el LLM:
%   - Tarea 1: sugerencias de reformulación de RA
%     → métricas automáticas (BLEU/ROUGE) + evaluación humana (Likert)
%   - Tarea 2: detección de inconsistencias RA ↔ contenidos
%     → precisión/recall sobre casos conocidos
%   - Tarea 3: generación de fundamentaciones sugeridas
%     → evaluación cualitativa de pertinencia y adecuación normativa
%   - Ejemplos representativos de buenas y malas sugerencias
%   - Análisis de alucinaciones o errores normativos del LLM
\ldots

%% ─────────────────────────────────────────────────────────────
\section{Resultados de búsqueda semántica (RAG)}
\label{sec:resultados_busqueda}
%% ─────────────────────────────────────────────────────────────

% TODO: Evaluar el módulo RAG:
%   - Precisión@k y Recall@k para consultas representativas
%   - Ejemplos de consultas y resultados recuperados
%   - Comparación con búsqueda por palabras clave (baseline)
%   - Análisis de calidad del contexto inyectado al LLM
\ldots

%% ─────────────────────────────────────────────────────────────
\section{Evaluación de trazabilidad y auditoría MCP}
\label{sec:resultados_trazabilidad}
%% ─────────────────────────────────────────────────────────────

% TODO: Evaluar la cobertura del log de auditoría:
%   - ¿Qué porcentaje de las acciones relevantes quedan registradas?
%   - ¿El log es suficientemente granular para una auditoría CONEAU?
%   - Ejemplo de log completo de un ciclo de vida de propuesta
%     (desde importación hasta aprobación)
%   - Análisis de overhead de rendimiento introducido por MCP
\ldots

%% ─────────────────────────────────────────────────────────────
\section{Evaluación de usabilidad del sistema}
\label{sec:resultados_usabilidad}
%% ─────────────────────────────────────────────────────────────

% TODO: Reportar evaluación de usabilidad (si se realizó):
%   - Perfil de los usuarios evaluadores (docentes, directores)
%   - Metodología: SUS, entrevistas, think-aloud, cuestionario propio
%   - Resultados: puntuación SUS o equivalente
%   - Tareas evaluadas y tasas de éxito
%   - Feedback cualitativo más relevante
%   - Si no se realizó evaluación formal: mencionar como limitación
%     y trabajo futuro
\ldots

%% ─────────────────────────────────────────────────────────────
\section{Análisis integrado y discusión}
\label{sec:discusion}
%% ─────────────────────────────────────────────────────────────

% TODO: Sintetizar y discutir los resultados:
%   - ¿Se respondieron las preguntas de investigación? (referenciar §3.2)
%   - ¿Qué dimensiones funcionaron mejor/peor y por qué?
%   - Comparación con trabajos relacionados de §2.4 y §2.5
%     (¿qué mejora esta solución respecto al estado del arte?)
%   - Limitaciones evidenciadas en los resultados:
%     * Dependencia del formato del template DOCX
%     * Latencia de llamadas al LLM
%     * Cobertura del corpus de evaluación
%   - Validez externa: ¿son generalizables los resultados a otras
%     carreras o instituciones?
\ldots

%% ─────────────────────────────────────────────────────────────
\section{Amenazas a la validez}
\label{sec:amenazas_validez}
%% ─────────────────────────────────────────────────────────────

% TODO: Describir amenazas a la validez del estudio:
%   - Validez interna: sesgos en la anotación del ground truth;
%     sobreajuste a las propuestas de \programaTesis
%   - Validez de constructo: ¿las métricas miden lo que queremos medir?
%   - Validez externa: corpus pequeño; un solo programa académico;
%     una sola institución
%   - Validez de conclusión: tamaño de muestra para métricas estadísticamente
%     significativas
\ldots
